\documentclass[conference]{IEEEtran}
\IEEEoverridecommandlockouts

\usepackage{cite}
\usepackage{amsmath,amssymb}
\usepackage{graphicx}
\usepackage{url}
\usepackage[hidelinks]{hyperref}
\usepackage{booktabs}
\usepackage{multirow}
\usepackage{algorithm}
\usepackage{algpseudocode}
\usepackage{xcolor}
\usepackage{array}

\newcommand{\phase}[2]{\par\smallskip\noindent\textbf{Phase #1: #2.}\par}
\newcommand{\step}[2]{\par\noindent\textit{Step #1: #2.}\ }

\begin{document}

\title{Context-Aware Alert Prioritization for Security Information and Event Management Systems}
\author{%
\IEEEauthorblockN{Anawin Suwannapha, Peeravich Teerapatanapan, Worathun Termjitcharoen,
Thanayos Ongpisut, Somchart Fugkeaw}
\IEEEauthorblockA{\textit{School of Information, Computer and Communication Technology (ICT)}\\
\textit{Sirindhorn International Institute of Technology, Thammasat University, Thailand}\\
6622780573@g.siit.tu.ac.th, 6622780649@g.siit.tu.ac.th, 6622790168@g.siit.tu.ac.th,\\
6622780672@g.siit.tu.ac.th, somchart@siit.tu.ac.th}
}

\maketitle

\begin{abstract}
Security Information and Event Management (SIEM) systems generate far more alerts than analysts can investigate, and ranking them by static rule severity does not reflect the operational risk of the underlying activity. Enrichment approaches attach attributes to each alert independently and so cannot distinguish an isolated event from the same event occurring within an attack. We present a context-aware prioritization layer for the open-source SIEM Wazuh that derives context across time within the alert stream itself. For each alert and entity, the layer forms a bounded window of the entity's recent activity, with repeats of one rule collapsed into a single entry, encodes it as a fifteen-feature vector computed from Wazuh alert fields alone, and scores it with an Isolation Forest trained without labels on benign activity and calibrated per host. The score is combined with a rule-novelty term, rule severity and asset exposure in a transparent hybrid function; alerts are presented as repeat groups so that floods occupy one queue position, and a per-entity accumulator escalates sustained activity. We evaluate on the public, labelled AIT Alert Data Set, eight simulated enterprise networks with multi-step attacks, under a tuning/test split fixed in advance, and characterize false promotion on five weeks of alerts from a real workstation.
\end{abstract}

\begin{IEEEkeywords}
alert prioritization, Wazuh, SIEM, alert correlation, anomaly detection, alert fatigue, MITRE ATT\&CK
\end{IEEEkeywords}

%% ======================================================================
\section{Introduction}

Security Information and Event Management (SIEM) systems are widely used in modern Security Operations Centers (SOCs) to collect, correlate, and analyze security events, but large alert volumes can overwhelm analysts and delay the investigation of high-risk incidents \cite{ref14}. A major limitation is that alerts are often ranked mainly by static rule severity or identifiers, which may not reflect the real operational risk of an event when factors such as asset criticality, user behavior, vulnerability exposure, timing, and surrounding attack patterns are considered. Prior research has therefore explored context-aware alert classification, adaptive SIEM correlation, contextual log filtering, reinforcement learning-based prioritization, and explainable graph-based triage to reduce alert fatigue and help analysts focus on alerts with greater security impact \cite{ref1,ref5,ref6,ref7,ref8}.

Wazuh provides a suitable base for this research: it supports log analysis, endpoint monitoring, rule-based detection, vulnerability detection and threat-intelligence integration. Recent Wazuh studies have integrated real-time threat feeds \cite{ref3}, evaluated coverage against MITRE ATT\&CK \cite{ref12} and applied machine learning to detection \cite{ref13}, but they target detection accuracy, enrichment or response automation rather than deciding which of the raised alerts should be investigated first.

A prioritization layer that evaluates each alert independently cannot express the distinction that matters most to an analyst. Consider two alerts identical in every recorded field: Wazuh rule 60122 (a failed Windows logon), the same host, the same user, the same severity level. In the first case it is the only such alert for that user in six hours, a mistyped password. In the second it is the fortieth within ninety seconds, followed by rule 60118, a successful logon. The second is a successful password-guessing attack and the first is routine noise, yet no attribute enrichment separates them, because the difference does not reside in the alert. It resides in the sequence of alerts surrounding it. Existing approaches therefore share three limitations:
\begin{enumerate}
\item \textbf{Context as attributes, not sequence.} Enrichment attaches asset, user or threat-intelligence attributes to each alert, but cannot see the alerts that precede it \cite{ref1,ref7}.
\item \textbf{Authored correlation.} Composite SIEM rules detect a multi-event pattern only if someone anticipated it and fixed its threshold in advance.
\item \textbf{Labels and volume.} Learned sequence models either need large event volumes and training budgets \cite{ref15}, a continuous expert feedback loop \cite{ref5}, or a purpose-built labelled corpus \cite{ref17}.
\end{enumerate}

This paper proposes a context-aware alert prioritization layer for Wazuh whose primary signal is derived across time within a single alert stream. The layer introduces no detections; it re-orders alerts that Wazuh has already raised. Each alert's preceding activity on the same entity is encoded as a window, rated by an unsupervised anomaly model and combined with rule novelty, rule severity and asset exposure in an interpretable priority, since analysts must be able to understand and challenge a ranking \cite{ref9,ref10}; a decaying per-entity risk accumulator raises sustained activity to incident level. The main contributions are:
\begin{itemize}
\item \textbf{A replayable system model.} A Wazuh prioritization layer in which correlation is derived mechanically from the alert stream rather than authored in advance as rule thresholds, with alert, context and risk state separated so that every result can be reproduced by replaying an immutable alert log.
\item \textbf{A temporal context window over collapsed runs.} A fifteen-feature window computed entirely from Wazuh alert fields, built from entries that collapse repeats of one rule, with an active-time correction that separates a host being switched off from an entity being quiet.
\item \textbf{An inspectable scoring and queueing design.} A per-host calibrated anomaly score combined with novelty, rule severity and exposure while remaining inspectable term by term, together with a repeat-group queue and an accumulator designed against measured double counting.
\item \textbf{A pre-specified evaluation.} An evaluation on a public labelled alert corpus, the AIT Alert Data Set \cite{ref21}, complemented by five weeks of benign alerts from a real workstation, with an ablation isolating the contribution of temporal context and comparators that test whether the learned model earns its place.
\end{itemize}

The remainder of the paper is organized as follows. Section~\ref{sec:related} reviews related work. Section~\ref{sec:formulation} states the problem and the preliminaries. Section~\ref{sec:system} presents the proposed system: its system model and the seven phases of its process. Section~\ref{sec:eval} describes the experiment design, and Section~\ref{sec:threats} discusses threats to validity, scope and future work.

%% ======================================================================
\section{Related Work}
\label{sec:related}

\subsection{Rule-Based Correlation in SIEM Platforms}

Wazuh supports composite rules through options such as \texttt{frequency} with \texttt{timeframe}, \texttt{if\_matched\_sid}, \texttt{if\_group}, \texttt{same\_source\_ip} and \texttt{same\_user}, so that a pattern spanning several events can raise a higher-severity alert. The limitation is authorship, not capability: each such rule requires someone to anticipate a pattern and fix its threshold, so only what has been imagined is detected. The approach proposed here is complementary: it forms a context window mechanically for every alert and learns from observed data what is unusual within it.

\subsection{Context-Aware Alert Prioritization}

Cinque et al. \cite{ref7} demonstrate contextual filtering and prioritization of application logs for situational awareness, establishing that context external to an individual record improves the assessment of that record. Suarez-Tangil et al. \cite{ref6} frame SIEM correlation as an adaptive process, arguing that static correlation cannot accommodate multi-step attacks. Roy and Balde \cite{ref1} apply large language models to context-aware alert classification, separating organizational context from fixed SIEM rules. These works establish the value of context but largely treat it as attributes attached to an alert, rather than as structure recovered from the alert sequence.

\subsection{Sequence Context and Alert Correlation}

DEEPCASE \cite{ref15} establishes the mechanism this work adopts: for each security event it builds a window of preceding events from the same device, identifies the relevant ones with an attention-based encoder-decoder, and clusters the attention vectors so that operators triage groups rather than events, reporting a 90.53\% reduction in analyst workload. AlertPro \cite{ref5} combines context features from alert sequences with history features from investigated alerts and finds that context features outperform raw alert fields. Both confirm that the preceding-event window carries discriminative signal, but DEEPCASE needs event volumes and a training budget that a small testbed cannot supply, and AlertPro depends on a continuous expert feedback loop.

\subsection{Machine Learning Applied to Wazuh}

Chamkar et al. \cite{ref13} apply machine learning to improve threat detection in Wazuh, and Al Siam et al. \cite{ref3} integrate real-time threat feeds with Wazuh for automated malware detection and response. Winkler and Sharma \cite{ref12} evaluate Wazuh against the MITRE ATT\&CK framework using adversary emulation. These works target detection accuracy or enrichment rather than the ordering of the resulting alert queue.

Closest to the present work, Smart-SIEM \cite{ref17} builds a per-source-IP behavioural context vector from the last $N$ events within Wazuh and classifies web application attacks with a LightGBM and XGBoost cascade. Its ablation over $N \in \{3,\dots,35\}$ selects $N = 30$, and context features improve all eight algorithms tested by 0.25 to 0.35 $F_1$, leading the authors to conclude that the context vector rather than the algorithm is the main contribution; that finding motivates the ablation design used here. The present work differs in three respects: it ranks rather than classifies; its model is trained without labels, where Smart-SIEM required 46{,}454 labelled events with synthetic oversampling; and it operates on the full Wazuh alert stream rather than HTTP telemetry, with an incident-level accumulation stage.

Graph-based approaches have also been applied to Wazuh telemetry, for attack-path prediction from event correlation \cite{ref4} and graph neural networks with continual learning on EDR data \cite{ref11}; a graph component was in this project's initial scope and was dropped in favour of learned scoring of stream-local context.

\subsection{Risk-Based Prioritization and Human Feedback}

Malach et al. \cite{ref8} combine alert triage with explainable graph representations and Shapley-based attribution, reinforcing that analysts require reasons accompanying a ranking. Shamsi et al. \cite{ref9} present human-centric risk-based alert prioritization in Splunk, and operational risk-based alerting practice \cite{ref19} accumulates risk per entity and escalates when a threshold is crossed together with evidence from multiple distinct sources or techniques. Khayat et al. \cite{ref10} combine multiple alert features using ensembled deep learning and hybrid optimization. Jalalvand et al. \cite{ref18} adapt prioritization through learning to defer with analyst feedback, and their survey \cite{ref16} organizes prioritization approaches by the degree of human-AI teaming involved. Feedback-driven adaptation is treated here as future work.

Closest in evaluation, Uetz et al. \cite{ref26} give the first systematic evaluation of risk-based alerting as continuous alert prioritization. They formulate five risk hypotheses (rule level, accumulation of alerts on an entity, variety of alert types, rarity of alert types, and aperiodicity), combine them, and evaluate on eight alert datasets, where the best combinations reach a mean AUROC of 0.92 against 0.72 for rule level alone. Their rarity and accumulation hypotheses correspond to the novelty signal in this design's window and to its accumulator. The present work differs in three respects: context is learned per host from the alert sequence rather than fixed by hypothesis; alerts are ranked as a queue of repeat groups and judged by when each attack stage is first met, rather than by alert-level AUROC; and all eight AIT-ADS scenarios are used under a tuning/test split fixed in advance, where they use one.

\subsection{Labelled Alert Corpora}

Public intrusion datasets such as CICIDS2017 and UNSW-NB15 consist of network flow records, not SIEM alerts; SALAD \cite{ref30} recasts their rows as SOC alerts with triage decisions but remains flow-derived, without SIEM rule identities or a baseline period before an attack. The AIT Alert Data Set (AIT-ADS) \cite{ref21} consists of the alerts that Wazuh, Suricata and AMiner raise when the logs of the AIT Log Data Set \cite{ref22} are replayed through them: eight simulated enterprise networks, each with several days of simulated user activity and one multi-step attack, labelled per time interval and per event. Its ground truth is a record of the attack rather than an analyst's judgement: scripted attacker models execute the attacks, and labelling rules filled in from each run's execution facts (attacker addresses and domains, start and stop times of each step) mark the log events each step caused \cite{ref22}, from which the alert labels follow. Uetz et al. \cite{ref26} report it as the only pre-existing dataset that met their requirements and use one scenario; Ndichu et al. \cite{ref27} derive a Wazuh alert stream from it. Microsoft's GUIDE \cite{ref28}, the largest public corpus of real incidents with analyst triage grades, since extended with priority labels \cite{ref29}, records Defender XDR incidents rather than Wazuh alerts and grades incidents rather than attack steps, so neither the rule term nor the per-host baseline applies to it. A public Wazuh alert set for fine-tuning language models \cite{ref31} spans under three hours with host identities removed, and its labels coincide with the rule level, so severity alone classifies it perfectly. Wazuh-collected Windows datasets \cite{ref23} label records by the presence of ATT\&CK tags on the triggering rule, which Section~\ref{sec:datasets} shows to be common in benign activity, so we do not use them as ground truth.

%% ======================================================================
\section{Problem Formulation and Preliminaries}
\label{sec:formulation}

\subsection{Alert Stream and Problem Statement}

Wazuh emits a time-ordered stream of alerts $a_1, a_2, \dots$. Each alert $a$ carries a rule identifier $r(a)$, a rule level $l(a) \in \{0,\dots,15\}$, rule groups $\mathrm{grp}(a)$, ATT\&CK techniques $\mathrm{tech}(a)$ and tactics $\mathrm{tac}(a)$, a timestamp $\tau(a)$, the reporting agent, and, when the decoder extracts them, a source address and a user. Given the stream up to time $t$, the layer must present an ordered queue of \emph{items}, repeat groups of alerts and risk notables, such that the items containing attack activity, weighted by what the attacker achieved, are met as early as possible. The layer re-orders alerts that Wazuh has raised; it neither creates detections nor reads log content. The design satisfies five requirements:
\begin{itemize}
\item \textbf{R1 Stream-local temporal context}: context is derived across time within the Wazuh alert stream; correlation with external log sources is out of scope.
\item \textbf{R2 No attack labels}: the scoring model is trained on benign baseline activity only.
\item \textbf{R3 Inspectability}: every priority decomposes into named terms that can be shown to the analyst.
\item \textbf{R4 Replayability}: all context and risk state is a pure function of the immutable alert log, so every system variant is produced by replaying one corpus.
\item \textbf{R5 Content-blindness}: only structured alert fields and detector rule identifiers are read, never command lines, paths or raw log lines.
\end{itemize}

\subsection{Isolation Forest}

The anomaly model is an Isolation Forest \cite{ref20}. A forest of $t$ trees is grown on subsamples of size $\psi$; each tree splits on a random feature at a uniformly random value until points are isolated, so a vector in a sparse region is isolated near the root. With $\ell(\mathbf{v})$ the path length of vector $\mathbf{v}$ and $c(\psi) = 2H(\psi - 1) - 2(\psi - 1)/\psi$ its expected value under unsuccessful search, $H$ the harmonic number,
\begin{equation}
s(\mathbf{v}) = 2^{-\mathbb{E}[\ell(\mathbf{v})] / c(\psi)} .
\label{eq:iforest}
\end{equation}
Unsupervised learning is used because it needs no attack labels and can rank activity no label described; it is the established choice for this role in alert triage \cite{ref5}.

%% ======================================================================
\section{Proposed System}
\label{sec:system}

This section presents the system model of the proposed prioritization layer and its operational process. Fig.~\ref{fig:architecture} shows the architecture.

\begin{figure}[t]
\centering
\includegraphics[width=0.92\columnwidth]{architecture.pdf}
\caption{Proposed prioritization layer. Tinted components retain state across alerts. The novelty term $S_\mathrm{nov}$ enters Component 7 from the same host window as $S_\mathrm{seq}$.}
\label{fig:architecture}
\end{figure}

\subsection{System Model}

The layer consists of nine components. Components 1--4 process each alert independently; Components 5, 6 and 8 retain state across alerts.
\begin{enumerate}
\item \textbf{Ingestion and Admission (C1)}: reads \texttt{alerts.json} from the Wazuh manager and discards alerts that do not describe behaviour.
\item \textbf{Normalization (C2)}: projects each admitted alert onto a fixed schema and assigns a total order.
\item \textbf{Entity Resolution (C3)}: maps each alert to the host, user and source-address entities it concerns.
\item \textbf{Entity Attribute Enrichment (C4)}: attaches history-independent attributes, giving the exposure score $S_\mathrm{exp}$.
\item \textbf{Temporal Context Window Builder (C5)}: maintains, per entity, a bounded window of recent activity over collapsed runs and encodes it as a feature vector.
\item \textbf{Sequence Anomaly Scoring (C6)}: scores each window with an Isolation Forest fitted on the baseline and calibrates the score per host, giving $S_\mathrm{seq}$.
\item \textbf{Hybrid Priority Scoring (C7)}: combines $S_\mathrm{rule}$, $S_\mathrm{exp}$, $S_\mathrm{seq}$ and $S_\mathrm{nov}$ into an inspectable priority.
\item \textbf{Entity Risk Accumulation (C8)}: keeps a decaying risk total per entity and raises risk notables.
\item \textbf{Repeat-Group Queue and Explanation (C9)}: presents repeat groups and notables to the analyst with their contributing terms.
\end{enumerate}

\textbf{State model.} The layer holds three kinds of state, summarized in Table~\ref{tab:state}. Context state and risk state are pure functions of the immutable alert log, the latter because the decay in Phase~7 is closed-form, so every system variant in the evaluation is produced by replaying one corpus under a different configuration (R4). Fitted models, calibration tables and tuned constants are \emph{parameters}: they are fitted once on baseline data in Phase~4, versioned, and never refitted on evaluation data. How each parameter value was chosen is listed in Table~\ref{tab:params} (Section~\ref{sec:params}). Table~\ref{tab:notation} summarizes the notation.

\begin{table}[t]
\caption{The three kinds of state held by the layer}
\label{tab:state}
\centering\footnotesize
\begin{tabular}{@{}llll@{}}
\toprule
& \textbf{Alert} & \textbf{Context} & \textbf{Risk} \\
\midrule
Key & alert id & (entity, alert) & entity \\
Behaviour & frozen & slides ($N$, $T$, $\delta$) & decays ($t_{1/2}$) \\
Produced by & C1--C4 & C5 & C8 \\
Scoring term & $S_\mathrm{rule}$, $S_\mathrm{exp}$ & $S_\mathrm{seq}$, $S_\mathrm{nov}$ & $R_e$ \\
\bottomrule
\end{tabular}
\end{table}

\begin{table}[t]
\caption{Notation used in the proposed system}
\label{tab:notation}
\centering\footnotesize
\setlength{\tabcolsep}{3pt}
\begin{tabular}{@{}l p{6.3cm}@{}}
\toprule
\textbf{Symbol} & \textbf{Description} \\
\midrule
$a$, $\tau(a)$, $l(a)$, $r(a)$ & alert, its timestamp, rule level and rule identity \\
$\mathcal{A}$ & admitted alert stream (Phase 1) \\
$E(a)$, $e$, $h(a)$ & entities of $a$, an entity, the host of $a$ \\
$\varepsilon$, $k$ & entry (collapsed run of one rule) and its member count \\
$\delta$, $N$, $T$ & entry gap, window length in entries, window expiry \\
$W(e,a)$, $\phi$ & window of $e$ at $a$, its feature vector \\
$G$, $I_h$, $g_i$ & downtime gap threshold, downtime intervals of host $h$, active gap between entry starts \\
$\mathcal{B}$ & (entity, rule) pairs seen in the baseline \\
$t$, $\psi$, $s(\mathbf{v})$ & trees, subsample size, raw forest score \\
$\hat F_{h,k}$, $\hat F_k$ & held-out score distributions of host $h$ and of entity type $k$ \\
$n_h$, $n_0$, $w_h$ & held-out windows of $h$, credibility constant, credibility weight \\
$S_\mathrm{seq}$, $S_\mathrm{nov}$ & calibrated sequence score, novelty term \\
$S_\mathrm{rule}$, $S_\mathrm{exp}$ & rule-level term, exposure term \\
$P(a)$, $w_r,w_c,w_s,w_n$ & hybrid priority and its weights \\
$\Delta$, $g$, $k_g$ & repeat-group gap, a group, its size \\
$\pi(g)$, $\lambda$, $\nu(g)$ & size-corrected group score, its saturation, group novelty flag \\
$R_e(t)$, $t_{1/2}$ & risk of entity $e$ at time $t$, decay half-life \\
$a_r, a_c, a_s$ & accumulator weights \\
$q(a,e)$, $n_e$ & window surprise, effective independent windows \\
$\theta$ & escalation threshold per entity type \\
\bottomrule
\end{tabular}
\end{table}

\subsection{System Process}

The layer operates in seven phases (Table~\ref{tab:phases}). Phases 1--3 apply identically to baseline and live alerts. Phase 4 is executed once per corpus on its baseline and produces the frozen parameters. Phases 5--7 run online for every admitted alert, using those parameters.

\begin{table}[t]
\caption{Phases of the process and the components that realize them}
\label{tab:phases}
\centering\footnotesize
\setlength{\tabcolsep}{3pt}
\begin{tabular}{@{}cllc@{}}
\toprule
\textbf{Phase} & \textbf{Purpose} & \textbf{Comp.} & \textbf{Mode} \\
\midrule
1 & Admission and normalization & C1, C2 & per alert \\
2 & Entity resolution and enrichment & C3, C4 & per alert \\
3 & Temporal context window construction & C5 & stateful \\
4 & Baseline fitting and calibration & C6 & offline, once \\
5 & Sequence scoring and hybrid priority & C6, C7 & online \\
6 & Repeat-group queue construction & C9 & online \\
7 & Entity risk accumulation and escalation & C8, C9 & stateful \\
\bottomrule
\end{tabular}
\end{table}

%% ---------------- Phase 1
\phase{1}{Alert Admission and Normalization (C1, C2)}
Not every Wazuh alert describes behaviour, and admitting those that do not corrupts the window features rather than merely diluting them. Phase 1 decides which alerts enter the layer and brings them into one schema.

\step{1}{Admission filter} The admitted stream is
\begin{multline}
\mathcal{A} = \{\, a : \mathrm{grp}(a) \cap (G_c \cup G_v) = \emptyset, \\
\mathrm{mode}(a) \neq \texttt{scheduled},\ a \notin \mathcal{D} \,\},
\label{eq:admission}
\end{multline}
where $G_c$ = \{\texttt{sca}, \texttt{rootcheck}, \texttt{policy\_monitoring}, \texttt{ciscat}\} are configuration-assessment groups, whose ATT\&CK fields describe what a check mitigates rather than what anyone did; $G_v$ = \{\texttt{vulnerability-detector}\} holds vulnerability scan results, which likewise describe an exploitable state rather than activity and on the workstation arrive in bursts (958 alerts in ten distinct minutes, levels up to 13) that would dominate a severity-ranked queue; none occur in the AIT-ADS tuning scenarios, and excluding them left every workstation measure unchanged; $\mathrm{mode}(a)$ is the file-integrity monitoring mode, excluded when \texttt{scheduled} because such an alert carries the time the scan ran rather than the time of the change (on the workstation, 1{,}190 registry alerts fell into four distinct minutes); and $\mathcal{D}$ is the set of duplicate ingestions, alerts that re-enter the stream through a second log path. In AIT-ADS every Suricata alert is ingested twice, once from \texttt{eve.json} and once from \texttt{fast.log}; the second copy is in $\mathcal{D}$. Excluded alerts are not discarded entirely: they still update the host's activity record used by the active-time correction of Phase 3.

\step{2}{Schema normalization} Each admitted alert is projected onto a fixed schema (rule identifier, level and groups; ATT\&CK techniques and tactics; timestamp; agent; source address; user) with an ingestion sequence number that totally orders alerts sharing a timestamp. ATT\&CK tags arrive as a nested object with tactic names from detection rules and as flat identifier lists from configuration checks; both are normalized to identifiers.

\step{3}{Rule identity of network alerts} Every Suricata alert reaches the layer under a single Wazuh rule (86601), so rule-based features cannot tell one network signature from another. The layer therefore uses the Suricata signature identifier as the rule identity $r(a)$ of these alerts (in the shaw scenario, 16{,}692 Suricata alerts carry ten signatures). It is a detector rule identifier, not log content, so the layer remains content-blind (R5).

%% ---------------- Phase 2
\phase{2}{Entity Resolution and Attribute Enrichment (C3, C4)}
Phase 2 determines the units over which context accumulates and attaches the attributes that do not depend on history.

\step{1}{Entity resolution} Each alert resolves to a set of entities $E(a) \subseteq \{\mathrm{host}, \mathrm{user}, \mathrm{srcip}\}$. The host is the monitored machine, mapped from the Wazuh agent through a frozen alias table so that a re-enrolled agent remains one entity. The user is taken, for Windows logon events, from the target account and otherwise from the subject account, with machine and service identities excluded and domain and case variants normalized. The source address is taken from the decoder's source field when routable. Table~\ref{tab:resolution} gives the measured resolution rates. Each corpus populates two of the three types, and the pair differs: the workstation has no network logons, and AIT's dominant benign alert, a mail-server login, records the account only in the raw line, which the layer does not read.

\begin{table}[t]
\caption{Measured entity resolution (share of admitted alerts)}
\label{tab:resolution}
\centering\footnotesize
\begin{tabular}{@{}lccc@{}}
\toprule
\textbf{Corpus} & \textbf{host} & \textbf{user} & \textbf{srcip} \\
\midrule
Workstation (25 days) & 100\% & 57.6\% & 0\% \\
AIT-ADS (two scenarios checked) & 100\% & 0.2--0.4\% & 52.7--63.1\% \\
\bottomrule
\end{tabular}
\end{table}

\step{2}{Exposure score} Component 4 gives an exposure score $S_\mathrm{exp} \in [0,100]$ from asset role. In AIT-ADS the dataset authors name the host behind every agent and place each role in a network zone \cite{ref22}; we set
\begin{equation}
S_\mathrm{exp} = 100\,(c + x)/2,
\label{eq:exposure}
\end{equation}
with exposure $x = 1$ for hosts reachable from the Internet (DMZ servers, firewall), $0.5$ for the intranet, reachable from outside only through the VPN, and $0$ for hosts the organization does not own, and criticality $c = 1$ for hosts that hold organizational data or control access, $0.5$ for the web server and $0$ for hosts not owned. The table was fixed from the role descriptions before any ranking was computed with it. On the single-host workstation $S_\mathrm{exp}$ is constant.

%% ---------------- Phase 3
\phase{3}{Temporal Context Window Construction (C5)}
For an admitted alert $a$ and an entity $e \in E(a)$, the window could simply be the last $N$ alerts concerning $e$, as in DEEPCASE \cite{ref15} and Smart-SIEM \cite{ref17}. On a busy host that window is a few seconds of one repeated rule: with $N = 20$, 89\% of baseline host windows in the four tuning scenarios held a single rule, and one flood evicts everything that preceded it. The window is therefore built from \emph{entries} rather than alerts. Algorithm~\ref{alg:window} gives the online procedure.

\step{1}{Run collapsing into entries} A \emph{run} of rule $r$ on entity $e$ is a maximal sequence of $e$'s admitted alerts with rule $r$ in which consecutive members are at most $\delta$ apart. An entry summarizes one run as
\begin{equation}
\varepsilon = \big(r,\ \max_{x} l(x),\ \textstyle\bigcup_x \mathrm{tech}(x),\ \bigcup_x \mathrm{tac}(x),\ \tau^{\mathrm{f}},\ \tau^{\mathrm{l}},\ k\big),
\label{eq:entry}
\end{equation}
its rule, highest level, ATT\&CK tags, first and last member times and member count. A new alert joins the most recent entry of its rule if $\tau(a) - \tau^{\mathrm{l}} \le \delta$, and otherwise opens a new entry; interleaved rules keep separate runs, and a run broken by a gap longer than $\delta$ is never rejoined.

\step{2}{Window assembly} The window $W(e,a)$ is the last $N$ entries of $e$, ordered by $\tau^{\mathrm{f}}$, whose newest member satisfies $\tau(a) - \tau^{\mathrm{l}} < T$, and it contains the entry holding $a$. We use $N = 20$, $\delta = 60$\,s and $T = 24$\,h (selected on the tuning scenarios, Section~\ref{sec:params}). An alert resolving to two entities updates two windows. A window is \emph{scoreable} when the entries it holds represent at least three alerts.

The entry gap $\delta$ is distinct from the queue's group gap $\Delta$ (Phase 6): groups merge over 24\,h for display, whereas entries keep $\delta$ short because their start times carry the timing features; at $N = 20$, raising $\delta$ from 60\,s to 900\,s lowered the mean AUC from 0.963 to 0.829.

\step{3}{Active-time correction} A workstation is switched off most nights, and a window opened by the first alert of the morning would otherwise span the overnight gap (up to 23\,h); the host being off is not the entity being quiet. For each host $h$ the builder records, online, every gap longer than $G$ in $h$'s \emph{global} stream (all its alerts, including those excluded at admission, which still prove the machine was on), giving downtime intervals $I_h$; a gap is recognized when the first alert after it arrives, so the correction is causal. For consecutive entries $\varepsilon_i, \varepsilon_{i+1}$ of the window, with start times $\tau_i = \tau^{\mathrm{f}}(\varepsilon_i)$,
\begin{equation}
g_i = \tau_{i+1} - \tau_i - \sum_{[s,f] \in I_h} \big| [\tau_i, \tau_{i+1}] \cap [s,f] \big| ,
\label{eq:activegap}
\end{equation}
the active gap between the starts of successive runs. Gaps lying wholly within downtime ($g_i = 0$ with downtime present) are omitted from the inter-arrival statistics, since they are not observations of the entity; counted as zero they would mimic two simultaneous alerts. $G = 1800$\,s was measured on the workstation, where it separated 23 nightly shutdowns from awake activity. The correction exists for the workstation only: it is disabled for AIT-ADS, whose servers run continuously, so no AIT result depends on $G$.

\step{4}{Feature extraction} Table~\ref{tab:features} defines the fifteen features over entries. Volume, which collapsing removes from the window's length, is kept explicitly in features 4 and 13. Features 2, 11, 12 and 13 are $\log(1+x)$-transformed before modelling. The tracked tactic ordinal follows ATT\&CK's Enterprise matrix order, and \texttt{tactic\_progression} is a weak ordering signal only, since tactics are not strictly sequential. Feature 15 is computed for display and is not a model input, so the model sees fourteen features.

\begin{table*}[t]
\caption{Window features over entries. $W = (\varepsilon_1,\dots,\varepsilon_n)$ ordered by start, $l_i$ the entry's highest level, $r_i$ its rule, $k_i$ its member count, $g_i$ the active gap between entry starts, $O$ the observed gaps, $\epsilon = 1$\,ms, $\varepsilon(a)$ the entry holding $a$}
\label{tab:features}
\centering\footnotesize
\begin{tabular}{@{}rlll@{}}
\toprule
\textbf{\#} & \textbf{Feature} & \textbf{Definition} & \textbf{Captures} \\
\midrule
1 & \texttt{window\_entries} & $n = |W|$ & sparse vs.\ saturated \\
2 & \texttt{window\_duration} & active time from $\tau^{\mathrm{f}}(\varepsilon_1)$ to $\tau(a)$ & burst vs.\ slow activity \\
3 & \texttt{distinct\_rule\_count} & $|\{r_i\}|$ & varied vs.\ repetitive \\
4 & \texttt{volume\_share\_max} & $\max_i k_i / \sum_i k_i$ & dominance of one run \\
5--6 & \texttt{rule\_level\_max}, \texttt{\_mean} & $\max_i l_i$,\ $\frac{1}{n}\sum_i l_i$ & severity \\
7 & \texttt{rule\_level\_trend} & OLS slope $\sum_i (i-\bar\imath)(l_i-\bar l) / \sum_i (i-\bar\imath)^2$ & escalation \\
8--9 & \texttt{distinct\_technique}/\texttt{tactic\_count} & $|\bigcup_i \mathrm{tech}(\varepsilon_i)|$,\ $|\bigcup_i \mathrm{tac}(\varepsilon_i)|$ & ATT\&CK breadth \\
10 & \texttt{tactic\_progression} & Spearman $\rho$(position, min tactic ordinal), tagged $\varepsilon_i$; 0 if $<3$ & kill-chain order \\
11--12 & \texttt{inter\_arrival\_mean}, \texttt{\_min} & mean and min of $\max(g_i,\epsilon)$ over $i \in O$ & pace, automation \\
13 & \texttt{current\_run\_length} & $k$ of $\varepsilon(a)$ & flood vs.\ isolated alert \\
14 & \texttt{novel\_rule\_flag} & $\mathbf{1}[(e, r(a)) \notin \mathcal{B}]$, $\mathcal{B}$ the (entity, rule) pairs seen in the baseline & novelty \\
15 & \texttt{night\_fraction} & share of entry starts outside 09:00--18:00 site time; \emph{not} a model input & --- \\
\bottomrule
\end{tabular}
\end{table*}

\begin{algorithm}[t]
\caption{Online window construction for one alert (Phase 3)}
\label{alg:window}
\begin{algorithmic}[1]
\Require alert $a$; per-entity deques $Q$; per-host last time $L$ and downtime $I$
\State $h \gets \mathrm{host}(a)$ \Comment{runs for every alert, admitted or not}
\If{$\tau(a) - L[h] > G$} \State $I[h] \gets I[h] \cup \{[L[h], \tau(a)]\}$ \EndIf
\State $L[h] \gets \max(L[h], \tau(a))$; drop intervals in $I[h]$ ending before $\tau(a) - T$
\If{$a \notin \mathcal{A}$} \Return \EndIf
\For{$e \in E(a)$}
  \State $\varepsilon \gets$ newest entry in $Q[e]$ with rule $r(a)$
  \If{$\varepsilon$ exists and $\tau(a) - \tau^{\mathrm{l}}(\varepsilon) \le \delta$} \State absorb $a$ into $\varepsilon$
  \Else \State append a new entry for $a$ to $Q[e]$; evict the oldest while $|Q[e]| > N$ \EndIf
  \State drop entries with $\tau(a) - \tau^{\mathrm{l}} \ge T$
  \State $\phi \gets \textsc{Features}(Q[e], a, I[h])$ \Comment{Table~\ref{tab:features}}
  \State emit $(a, e, \phi)$, scoreable iff $\sum_{\varepsilon \in Q[e]} k(\varepsilon) \ge 3$
\EndFor
\end{algorithmic}
\end{algorithm}

%% ---------------- Phase 4
\phase{4}{Baseline Fitting and Calibration (C6, offline)}
Phase 4 runs Phases 1--3 over a corpus's baseline, the activity before its first attack step in AIT-ADS and the frozen baseline on the workstation, and fits every quantity the online phases need. Its output is versioned and frozen; nothing in it is refitted on evaluation data. Algorithm~\ref{alg:fit} summarizes it.

\step{1}{Training-set construction} One forest is fitted per entity type, since dense host windows and sparse address or user windows follow different distributions. Successive windows of an entity share almost all their entries, so training uses the scoreable baseline windows at a stride of 10; the windows the stride leaves out are \emph{held out} for calibration (Step 3). Where several hosts feed one forest, each contributes at most an equal share of the strided windows, taken evenly over its baseline: unstratified, the three busiest hosts of the AIT tuning scenarios would supply 64--89\% of host windows and the intranet server, the file share and the cloud share under 1\% each, so the mail server and the firewall would define normal for every host. Inputs are the fourteen model features of Table~\ref{tab:features}, unscaled. A forest is fitted only for an entity type with at least 200 training windows, which excludes AIT's user entities (5--26).

\step{2}{Forest fitting} Each forest uses $t = 200$ trees and $\psi = \min(4096, n)$ for $n$ training windows. The usual $\psi = 256$ assumes anomalies form a dense cluster, which attack windows do not; in simulation over ten seeds it ranked attacks at AUC 0.926, below a Mahalanobis comparator (0.965), while $\psi = 4096$ reached 0.991. On the workstation the training target is 8{,}000 windows after striding, with a floor of 4{,}000 (in simulation, below about 2{,}000 the seed variance exceeded the effect the ablation measures). An AIT scenario's pre-attack period is fixed: the tuning scenarios give 1{,}840--6{,}885 host and 583--3{,}874 source-address windows after striding, mostly below the floor, so the spread across seeds is reported rather than assumed (up to 0.031 in $\alpha$-nDCG@20 on the tuning set). Each seed is an independent replica, forest and calibration both; metrics are computed per seed and reported as mean $\pm$ standard deviation over ten seeds, never on seed-averaged scores, which would amount to a larger forest and hide the variance.

\step{3}{Per-host credibility calibration} Raw scores concentrate near 0.5, so a linear rescale would make the sequence term nearly constant. Scores are instead mapped through empirical distributions of the held-out baseline windows, with the host's own distribution shrunk towards the pooled one in proportion to how little evidence supports it:
\begin{equation}
S_\mathrm{seq}(a, e) = 100 \big[ w_h \hat F_{h,k}(s) + (1 - w_h) \hat F_{k}(s) \big],
\label{eq:calib}
\end{equation}
where $s = s(\phi(W(e,a)))$ by (\ref{eq:iforest}), $\hat F_{h,k}$ is fitted on the $n_h$ held-out windows of host $h$ and entity type $k$, $\hat F_k$ on those of all hosts of type $k$, and $w_h = n_h/(n_h + n_0)$, with $n_0 = 100$ the number of windows at which the two weigh equally. The weight is B\"uhlmann's credibility weight \cite{ref25}: a host's own estimate is trusted in proportion to its sample rather than all or nothing. Calibrating on the training windows themselves would be biased, since the trees were grown from those points; the held-out windows remove that bias at no cost, so the workstation's drift check measures behaviour rather than an artefact.

For a host with a long baseline $w_h \to 1$, $S_\mathrm{seq}$ is uniform on its benign data and expresses how unusual a window is \emph{for that machine}. For a host with few windows the pooled distribution dominates, and the score largely states that the host is rarely active. In AIT-ADS this applies to the intranet server, the file share and the cloud share, with 17--116 held-out host windows each ($w_h = 0.15$--$0.54$) across the tuning scenarios. Shrinkage replaced a hard switch (the host's own distribution from 500 held-out windows, the pooled one below), under which those hosts' benign evaluation windows exceeded the 99th percentile in 36--100\% of cases instead of about 1\%; it lowers this to 0--50\% wherever a host has more than four benign evaluation windows, reducing the inflation without removing it, and where baselines are short calibration also drifts (4--42\% in russellmitchell even for hosts with long baselines). Results are therefore reported separately for stages on hosts calibrated mainly by their own baseline ($n_h \ge n_0$) and on the rest (Section~\ref{sec:stats}).

\step{4}{Baseline reference sets} Phase 4 also records the set $\mathcal{B}$ of (entity, rule) pairs seen in the baseline, used by feature 14 and the novelty term; and, after replaying the baseline through Phases 5--7, the per-entity-type distribution of the accumulated risk $R_e(t)$, whose 99.9th percentile is the escalation threshold $\theta$ and which places risk notables on the queue's scale (Phase 7).

\begin{algorithm}[t]
\caption{Baseline fitting for one corpus and seed (Phase 4)}
\label{alg:fit}
\begin{algorithmic}[1]
\Require baseline alerts; seed; $t$, $n_0$
\Ensure forests $\{f_k\}$, calibration tables, $\mathcal{B}$, $\{\theta_k\}$
\State run Phases 1--3 over the baseline; keep scoreable windows
\State $\mathcal{B} \gets \{(e, r(a))\}$ over admitted baseline alerts
\For{each entity type $k$}
  \State $X_k \gets$ every 10th window of each entity; $H_k \gets$ the rest
  \State cap each host's share of $X_k$ at an equal share, evenly over time
  \If{$|X_k| < 200$} mark type $k$ unscored; \textbf{continue} \EndIf
  \State $f_k \gets \textsc{IsolationForest}(X_k, t, \psi = \min(4096, |X_k|), \text{seed})$
  \State $\hat F_k \gets$ ECDF of $s_{f_k}$ over $H_k$
  \For{each host $h$}
    \State $\hat F_{h,k} \gets$ ECDF over $h$'s windows in $H_k$; $w_h \gets n_h/(n_h + n_0)$
  \EndFor
\EndFor
\State replay the baseline through Phases 5--7
\State $\theta_k \gets$ 99.9th percentile of $R_e(t)$ over entities of type $k$
\State freeze and version all outputs
\end{algorithmic}
\end{algorithm}

%% ---------------- Phase 5
\phase{5}{Sequence Scoring and Hybrid Priority (C6, C7)}
Phase 5 turns the windows of an admitted alert into calibrated scores and combines them with the per-alert attributes into one priority.

\step{1}{Window scoring} Every scoreable window $W(e,a)$ is scored with the forest of $e$'s type by (\ref{eq:iforest}) and calibrated by (\ref{eq:calib}), giving $S_\mathrm{seq}(a,e) \in [0,100]$.

\step{2}{Reduction to the host window and cold start} Every entity's window is scored, but the score entering the queue is that of the host window, $S_\mathrm{seq}(a) = S_\mathrm{seq}(a, h(a))$; source-address and user windows feed only the accumulator (Phase 7), where following an external address across hosts is their purpose. The originally proposed maximum over the alert's entities was dropped on the tuning scenarios: source-address windows alone rank the non-scan stages no better than rule severity ($\alpha$-nDCG@20 0.165 against 0.198) and inside the maximum lift benign groups above the attack, whereas host windows alone placed the first item of six of nine non-scan stages earlier and none later ($p = 0.031$). Windows representing fewer than three alerts receive the baseline median 50 and are flagged, so that the analyst sees ``no temporal evidence'' rather than ``normal''.

\step{3}{Hybrid priority} Component 7 combines
\begin{equation}
P(a) = w_r\, S_\mathrm{rule}(a) + w_c\, S_\mathrm{exp}(a) + w_s\, S_\mathrm{seq}(a) + w_n\, S_\mathrm{nov}(a),
\label{eq:hybrid}
\end{equation}
with $S_\mathrm{rule} = 100\, l(a)/15$, linear so that each term's contribution can be displayed, and $S_\mathrm{nov} = 100$ when the host window flags the alert's rule as never seen on that host in the baseline (feature 14), 0 otherwise. The novelty flag is stated as a term of its own because it is the most influential feature for 53 of the top 80 queue items on the tuning scenarios and on its own separates attack from benign alerts within a host better than $S_\mathrm{seq}$. The weights, selected on the tuning scenarios over the weight simplex (Section~\ref{sec:params}), are $(w_r, w_c, w_s, w_n) = (0, 0, 0.8, 0.2)$: rule severity never helped and no exposure weight passed the selection protocol, so on AIT-ADS the priority consists of the sequence and novelty terms. The rule and exposure terms stay in the form and are shown in every item's explanation, for inspectability (R3) and because a deployment with a real asset inventory may weight exposure differently; the default weights $(0.3, 0.2, 0.5, 0)$ are reported beside the tuned ones. The novelty term shares the host-rarity confound (Section~\ref{sec:threats}): on benign groups the flag fires two to thirty times as often on the rarely active hosts (12--57\% of their groups against 1--21\%), which is why the within-host AUC and the calibration split are reported.

\step{4}{Attribution} For an alert an analyst opens, each feature $j$ is replaced by its baseline median and the drop in $S_\mathrm{seq}$ is reported; the three largest drops are shown.

%% ---------------- Phase 6
\phase{6}{Repeat-Group Queue Construction (C9)}
An attack phase or a benign job can raise thousands of near-identical alerts; in AIT-ADS a single directory brute force produces over 400{,}000. Ranked individually they occupy the whole queue. Phase 6 therefore ranks \emph{repeat groups}.

\step{1}{Repeat grouping} A repeat group is a maximal run of admitted alerts with the same host and rule in which consecutive members are at most $\Delta = 24$\,h apart. Each group is shown once, with its count and time span. On the four tuning scenarios this reduces 13{,}877--466{,}383 alerts in the evaluation period to 51--85 queue items.

\step{2}{Size-corrected group score} A group's priority must summarize up to hundreds of thousands of member scores. The maximum gives a group of $k_g$ members $k_g$ chances to draw a high score; a group is therefore scored by the surprise of its maximum given its size,
\begin{equation}
\pi(g) = -\log_{10}\!\big(1 - u_g^{\,k_g}\big), \qquad u_g = \max_{a \in g} S_\mathrm{seq}(a)/100,
\label{eq:pi}
\end{equation}
how improbable it is that $k_g$ benign alerts would produce a maximum this high, left uncapped when ranking $S_\mathrm{seq}$ alone so that the top of the queue has no ties. This is the same correction Phase 7 applies to its window term. The choice is empirical: $n = k_g$ treats members as independent although their windows overlap, yet the better-motivated discount $n = k_g/N$ performed worse; and any size discount, like the maximum before it, can hide a single anomalous alert inside a large group of the same rule on the same host. The maximum itself is carried as the ablation arm B2$_\mathrm{max}$.

\step{3}{Group priority and explanation} Every member of a group shares its host and rule, so $S_\mathrm{rule}$ and $S_\mathrm{exp}$ are constant within it, and the group's priority is
\begin{multline}
P(g) = w_r S_\mathrm{rule}(g) + w_c S_\mathrm{exp}(g) \\
+ w_s \cdot 100 \min\!\big(1, \pi(g)/\lambda\big) + w_n \cdot 100\,\nu(g),
\label{eq:grouppriority}
\end{multline}
the sequence term on the same 0--100 scale as $q$ in Phase 7, with $\lambda = 8$ and $\nu(g) = 1$ when a member's host window flags its rule as novel. Each queue item lists its three highest-contributing terms, the window features that most influenced $S_\mathrm{seq}$, and, for risk notables, the contributing groups.

%% ---------------- Phase 7
\phase{7}{Entity Risk Accumulation and Escalation (C8)}
Phase 7 keeps a decaying risk total per entity so that a campaign of individually unremarkable alerts becomes visible. Algorithm~\ref{alg:accum} gives the per-alert update.

\step{1}{Double-counting remedies} Summing $P(a)$ over alerts, as first proposed, measures volume rather than evidence: on a 300-alert sample it inflated host risk twentyfold, 93\% of it from one 23-second configuration scan. Three remedies follow. \textit{Repeat inflation:} alerts are grouped into accumulation groups $g$ by (entity, rule), split when the rule is silent for longer than $T$, and a group of $k_g$ alerts contributes $(1 + \ln k_g)$ times one alert's value. \textit{Window evidence charged per alert:} a window persists across many alerts, so the window term enters as a decayed running maximum, not a sum. \textit{Cross-entity fan-out:} totals are compared only within an entity type, and notables sharing contributing alerts are merged into one incident.

\step{2}{Accumulated risk} The accumulator is
\begin{multline}
R_e(t) = \sum_{g \in \mathcal{G}_e} v_g\,(1 + \ln k_g)\, 2^{-(t - t^{\mathrm{last}}_g)/t_{1/2}} \\
+ a_s \max_{a :\, t - \tau(a) < T} q(a, e)\, 2^{-(t - \tau(a))/t_{1/2}},
\label{eq:accum}
\end{multline}
with $\mathcal{G}_e$ the accumulation groups of $e$, $v_g = a_r S_\mathrm{rule} + a_c S_\mathrm{exp}$ of the group's first member, $t^{\mathrm{last}}_g$ its most recent member, and half-life $t_{1/2} = 1$\,h. The accumulator keeps its own weights, $(a_r, a_c, a_s) = (0.3, 0.2, 0.5)$, separate from the queue's: the tuned queue weights give the static terms zero weight (Phase 5), which would leave the accumulator no group evidence at all. Decaying from the most recent member and splitting groups at gaps longer than $T$ keep a burst today from being absorbed into a group that began weeks earlier.

The window term uses a surprise score rather than $S_\mathrm{seq}$, whose maximum over many windows approaches 100 for any busy entity because $S_\mathrm{seq}$ is uniform on benign data. With $u = S_\mathrm{seq}(a,e)/100$ and $n_e$ the effective number of independent windows the entity produced in the preceding $T$,
\begin{equation}
q(a, e) = 100 \cdot \min\!\big(1,\ -\log_{10}(1 - u^{\,n_e}) / \lambda \big),
\label{eq:q}
\end{equation}
the improbability that $n_e$ benign windows would produce a maximum this high, on a log scale saturating at $\lambda = 4$. A window changes materially only when an entry is added to it, so $n_e$ is the number of entries the entity opened in the preceding $T$ divided by $N$ (for B2$_\mathrm{alert}$, the window count divided by $N$). This correction is derived rather than measured; it is validated by checking that the window term is not flat for the workstation host across a benign day.

\step{3}{Escalation} An entity becomes a risk notable when three conditions hold together: $R_e(t) > \theta$, an amount of evidence; its contributing groups number at least two, so that one flood cannot escalate on its own; and they carry at least two ATT\&CK tactics contributed by \emph{different} rules. The last condition is a correction: several common benign rules each carry two or three tactics, so a plain ``two tactics'' condition is met by a single benign alert. $\theta$ is calibrated as $S_\mathrm{seq}$ is, per entity type, as the 99.9th percentile of $R_e(t)$ over the baseline period, so that it states a benign crossing rate rather than a hand-set level. A group is \emph{contributing} only while its decayed term is at least 10\% of $R_e(t)$, so that a residue from the previous day cannot supply the second group.

\step{4}{Notables in the queue} A risk notable enters the queue as one item that absorbs its contributing groups, which are listed inside it rather than ranked separately, following risk-based alerting practice in which the incident, not its contributing events, is what the analyst triages \cite{ref19}. Its priority is placed on the queue's scale by the same device as $S_\mathrm{seq}$, the percentile of $R_e(t)$ among the entity type's baseline values. Since $\theta$ is the 99.9th percentile, a notable enters at 99.9 or above, in practice above nearly every group, and notables are ordered among themselves by that percentile.

\begin{algorithm}[t]
\caption{Risk accumulation and escalation for one admitted alert (Phase 7)}
\label{alg:accum}
\begin{algorithmic}[1]
\Require alert $a$; $S_\mathrm{seq}(a,e)$ for $e \in E(a)$; frozen $\theta_k$
\For{$e \in E(a)$}
  \State $g \gets$ accumulation group of $(e, r(a))$; open a new one if the rule was silent longer than $T$
  \State on opening: $v_g \gets a_r S_\mathrm{rule}(a) + a_c S_\mathrm{exp}(a)$; $k_g \gets 0$
  \State $k_g \gets k_g + 1$; $t^{\mathrm{last}}_g \gets \tau(a)$
  \State $n_e \gets$ (entries $e$ opened in the last $T$) $/ N$
  \State compute $q(a,e)$ by (\ref{eq:q}); retain $(\tau(a), q(a,e))$ for time $T$
  \State compute $R_e(\tau(a))$ by (\ref{eq:accum})
  \State $C \gets$ groups whose decayed term is $\ge 0.1\,R_e$
  \If{$R_e > \theta_{k(e)}$ \textbf{and} $|C| \ge 2$ \textbf{and} $C$ has $\ge 2$ tactics from different rules}
    \State raise or update notable for $e$, absorbing $C$
    \State merge notables sharing contributing alerts
    \State place it at the baseline percentile of $R_e$
  \EndIf
\EndFor
\end{algorithmic}
\end{algorithm}

%% ======================================================================
\section{Experiment Design}
\label{sec:eval}

The experiments are designed to answer five questions: (RQ1) does temporal context rank attack stages earlier than rankings without it, including host and level rarity; (RQ2) does the Isolation Forest earn its place over a simple robust z-score on the same windows; (RQ3) which design choices, entries, the size-corrected group score and the novelty term, are necessary; (RQ4) does entity risk accumulation surface incidents, and how does the method compare with a hand-designed risk-based alerting score; and (RQ5) how much benign activity on a real machine would the layer promote.

\subsection{Datasets}
\label{sec:datasets}

The two corpora have different roles and never meet: AIT-ADS supplies every ranking result, and the workstation, which contains no attacks, answers RQ5. Each is replayed separately through the seven phases with its own fitted models. Only configuration crosses, in one direction: the weights, $\Delta$ and $\theta$ selected on the AIT tuning scenarios are applied to the workstation unchanged, so that its false-promotion figures describe the layer as it would be deployed. Table~\ref{tab:corpora} lists every setting that differs.

\begin{table}[t]
\caption{Corpus roles and the settings that differ. All other components and parameters are identical}
\label{tab:corpora}
\centering\footnotesize
\setlength{\tabcolsep}{3pt}
\begin{tabular}{@{}lll@{}}
\toprule
& \textbf{AIT-ADS (primary)} & \textbf{Workstation (secondary)} \\
\midrule
Role & ranking evaluation & false promotion \\
Labels & event-level & none (no attacks) \\
Hosts & 9--11 Ubuntu servers, & 1 Windows 11, \\
 & simulated users & real user \\
\midrule
Ph.1 duplicates $\mathcal{D}$ & Suricata \texttt{fast.log} & none \\
Ph.1 time / site zone & \texttt{@timestamp} / UTC & \texttt{timestamp} / UTC+7 \\
Ph.2 alias table & none & agent 002 $\to$ 001 \\
Ph.2 entities populated & host, srcip & host, user \\
Ph.2 $S_\mathrm{exp}$ & documented host roles & constant (one host) \\
Ph.3 active time & off & on, $G = 1800$\,s \\
Ph.4 training data & pre-attack period & own baseline \\
Ph.5--7 $w$, $\Delta$, $\theta$ & selected (tuning set) & inherited, frozen \\
\midrule
Metrics & coverage@$k$, ranks, AUC & drift, load per awake hour \\
\bottomrule
\end{tabular}
\end{table}

\textbf{AIT-ADS (primary)} \cite{ref21} provides the labelled evaluation. Each of its eight scenarios is a separate simulated organization of 9--11 monitored hosts with 1.9--4.3 days of activity before a multi-step attack of ten steps: network and service scans, directory and WordPress scans, a webshell upload, password cracking, a reverse shell, privilege escalation, a service stop, and DNS exfiltration. Table~\ref{tab:ait} summarizes the scenarios after admission. We use the event-level labels, which mark the alerts produced by attacker actions, and not the time-interval labels, which mark every alert inside an attack phase: we found that most time-labelled alerts in the post-reconnaissance phases are benign mail logins and antivirus updates that happened to occur during the phase.

\textbf{What reaches the layer.} Which attack steps produce an alert at all is a property of the corpus, and we measured it before designing the metrics. Of the ten steps, four (the network scan, webshell upload, reverse shell and service stop) raise no Wazuh or Suricata alert in any of the eight scenarios. Password cracking raises one only where it is performed online against a service (two of the seven scenarios that include it, through Suricata), and the exfiltration in two of eight. The directory and WordPress scans and privilege escalation raise alerts in all eight, the service scan in seven. A layer that re-ranks alerts cannot rank what was never raised, so stage coverage counts only stages with at least one alert, and the unalerted steps are reported as outside the layer's boundary rather than as misses. The gap lies in the rule set, not in the logs: the dataset's log-based anomaly detector, AMiner, run on the same raw logs, flags the webshell commands, offline password cracking and DNS scanning that Wazuh does not.

\textbf{Designed contamination.} The dataset authors start the DNS exfiltration at the beginning of every simulation, deliberately poisoning the training period of anomaly detectors \cite{ref22}. It reaches the alert stream only in harrison and santos, where a Suricata rule for DNS queries to a \texttt{.biz} domain raised 46{,}211 and 25{,}707 event-labelled exfiltration alerts \emph{before} the first attack step, 52\% and 49\% of their baselines; these two scenarios test the layer under a heavily contaminated baseline. Since a deployed layer cannot know its baseline is clean, the recorded condition is the headline: the test run uses recorded baselines, and santos is run a second time with its labelled exfiltration alerts removed from training to measure what the contamination costs. Tuning used harrison oracle-cleaned, with its recorded run reported beside it. The end of the exfiltration, which the authors intend as a detection target, produces an absence of alerts and is outside what an alert-ranking layer can observe.

\begin{table}[t]
\caption{AIT-ADS scenarios after admission. Base: alerts before the first attack step ($^\ast$about half are exfiltration alerts). Causal: event-labelled attack alerts other than scans (exfiltration in parentheses, almost all before the first step)}
\label{tab:ait}
\centering\footnotesize
\setlength{\tabcolsep}{4.5pt}
\begin{tabular}{@{}lrrrrr@{}}
\toprule
\textbf{Scenario} & \textbf{Hosts} & \textbf{Span (d)} & \textbf{Admitted} & \textbf{Base} & \textbf{Causal} \\
\midrule
\multicolumn{6}{@{}l}{\textit{tuning set}}\\
harrison & 9 & 5.0 & 516{,}343 & 89{,}476$^\ast$ & 13 (47{,}183) \\
wheeler & 11 & 5.0 & 535{,}265 & 68{,}882 & 9 \\
russellmitchell & 10 & 4.0 & 32{,}302 & 18{,}425 & 6 \\
shaw & 11 & 6.0 & 51{,}831 & 28{,}682 & 10 \\
\midrule
\multicolumn{6}{@{}l}{\textit{test set}}\\
fox & 11 & 4.9 & 451{,}439 & 18{,}790 & 12 \\
wilson & 10 & 6.0 & 560{,}265 & 79{,}901 & 10 \\
santos & 10 & 3.9 & 82{,}401 & 52{,}121$^\ast$ & 14 (25{,}707) \\
wardbeck & 10 & 5.0 & 63{,}694 & 23{,}463 & 10 \\
\bottomrule
\end{tabular}
\end{table}

\textbf{Workstation (secondary).} A Windows 11 workstation in daily personal use was monitored with stock Wazuh 4.14.7 and Sysmon (alert level 3, no custom rules) from 12 September 2026, with the baseline frozen on 17 October. At the freeze it held 128{,}170 admitted alerts over 35 days (about 3{,}700 per calendar day, at about 10 awake hours per day; 127{,}212 once vulnerability-detector results are excluded in Phase~1). Its final week is held out (21{,}617 admitted alerts, 68.9 awake hours); the rest gives 106{,}551 host and 61{,}590 user scoreable windows, about 10{,}700 and 6{,}200 after striding. It contains no attacks and is used for what AIT-ADS cannot show: false promotion on a real person's activity. With one host, B1 and B0 rank identically there. Measured on it: 7.7\% of benign alerts carry ATT\&CK tags, the most frequent being T1570 \emph{Lateral Tool Transfer} on a machine with nothing to move laterally to, and 34\% of alerts arrive within fifteen minutes of the machine starting.

\subsection{Systems Compared}

\begin{itemize}
\item \textbf{B0}: rule level. \textbf{B1}: B0 + exposure.
\item \textbf{B0$_\mathrm{host}$}: groups ranked only by how rarely their host alerted in the baseline, $100\,(1 - \text{host's share of baseline alerts})$. It tests whether temporal context adds anything beyond ``this machine is unusual'': in AIT-ADS the attack lands on the least active host.
\item \textbf{B0$_\mathrm{rank}$}: B0 with the level replaced by its rarity, the percentile of the alert's level among its host's baseline alerts. It uses no temporal context; on the tuning set it and B0$_\mathrm{host}$ were the strongest systems without it, so B2 must beat them, not only B0.
\item \textbf{B2}: the method, the sequence and novelty terms with the tuned weights (Phase 5); \emph{B2 vs.\ B1 isolates temporal context}.
\item \textbf{B2$'$}: B2 with the forest replaced by a robust z-score, $\max_j |z_j|$ against baseline medians: does the forest earn its place?
\item \textbf{B2$_\mathrm{nov}$}: B2 with $S_\mathrm{seq}$ replaced by the novelty flag alone: does the window add anything beyond ``never seen this rule here''?
\item \textbf{B2$_\mathrm{max}$}: B2 with groups ranked by their maximum $S_\mathrm{seq}$ instead of $\pi(g)$: does the size correction earn its place?
\item \textbf{B2$_\mathrm{seq}$}: B2 without the novelty term, the sequence term alone: does the novelty term add to the window?
\item \textbf{B2$_\mathrm{alert}$}: B2 with the window of the last 5 alerts, the strongest alert window on the tuning set, instead of entries: does collapsing change the ranking?
\item \textbf{B3}: B2 + entity accumulation (Phase 7).
\item \textbf{B$_\mathrm{RBA}$}: risk-based alerting after Uetz et al. \cite{ref26}: their five hypotheses scored per alert as in their implementation and combined by a weighted mean whose windows and weights are selected on the tuning set, groups taking their members' maximum. For comparability its windows look only into the past (their better centred windows also use future alerts) and every normalizer is estimated from the pre-attack period. Does a learned window beat a strong hand-designed risk score?
\end{itemize}

\subsection{Evaluation Protocol}
\label{sec:protocol}

For each scenario, the forest and calibration are fitted on that scenario's activity before its first attack step and frozen (Phase 4); the evaluation period is the remainder. Scenarios are split before any tuning into a tuning set \{harrison, wheeler, russellmitchell, shaw\} and a test set \{fox, wilson, santos, wardbeck\}, each with two large and two small scenarios and one with exfiltration. All weights, $\Delta$, $\theta$ and the novelty and calibration choices are selected on the tuning set; every reported comparison comes from the test set. Ranking is measured over queue items. Ground truth comes from the event labels, graded by what the attacker's action achieved and fixed before any test result is seen (Table~\ref{tab:grades}): an item's grade is the highest grade among its members' stages, and benign items have grade 0. Every event label that occurs among the Wazuh alerts of the eight scenarios has a grade. The grades derive from the attack steps described by the dataset authors and from nothing the layer reads, so no system is favoured by construction, as one would be if items were graded by rule level or by the ATT\&CK tags of the triggering rule.

\begin{table}[t]
\caption{Relevance grades, fixed before testing}
\label{tab:grades}
\centering\footnotesize
\begin{tabular}{@{}c p{4.1cm} p{2.6cm}@{}}
\toprule
\textbf{Grade} & \textbf{Stages (event labels)} & \textbf{Attacker has} \\
\midrule
3 & user change, sudo escalation, exfiltration & control or data \\
2 & online password cracking & a credential attack \\
1 & service scan, directory scan, WordPress scan & reconnaissance only \\
0 & benign & --- \\
\bottomrule
\end{tabular}
\end{table}

\textbf{Selection objective.} The tuning objective, for weights, $\Delta$, $\theta$ and the remaining choices, is the mean graded $\alpha$-nDCG@20 over the tuning scenarios. A difference smaller than the larger of the two configurations' seed standard deviations counts as a tie, broken by MRR when the MRR difference exceeds its larger seed standard deviation, and otherwise resolved in favour of the incumbent. Because the primary test is the paired first-item rank (Section~\ref{sec:stats}), a configuration the objective prefers is vetoed if it places the first item of a non-scan stage later on a majority of stage occurrences than the incumbent. The veto was added on the tuning set, before any test run, when the objective preferred a group score that met eight of nine stages later.

\subsection{Metrics}

\begin{itemize}
\item \emph{Graded $\alpha$-nDCG@k}, $k \in \{10, 20, 50\}$, the primary metric. With $g_s$ the grade of stage $s$ and $c_{s,i}$ the number of items above rank $i$ that contain $s$, the gain at rank $i$ is $\sum_{s \in i} g_s (1-\alpha)^{c_{s,i}}$, discounted by $\log_2(i+1)$ and normalized by the greedily computed ideal ordering \cite{ref24}. The factor $(1-\alpha)^{c}$, with $\alpha = 0.5$, credits the first item of a stage fully and each repeat less, so a scan that recurs on five hosts does not count five times. An item that bundles several stages earns each stage's gain at its rank, so risk notables, which absorb groups, gain from bundling as well as from position; first-item ranks are reported beside $\alpha$-nDCG for that reason.
\item \emph{First-item rank} of each non-scan stage and their mean reciprocal rank (MRR): how far the analyst scrolls before first meeting each stage that matters.
\item \emph{Stage coverage@k}, reported separately for non-scan and scan stages, since scans are surfaced by every system and would otherwise dominate.
\item \emph{Macro-AUC}: the alert-level AUC of each non-scan stage against benign alerts, averaged over stages so that a stage with thousands of alerts does not outweigh one with ten. It is queue-independent and secondary.
\item \emph{Within-host AUC}: for each host carrying at least five non-scan attack alerts and five benign alerts, the alert-level AUC of its non-scan attack alerts against its own benign alerts, averaged over those hosts. Host rarity is constant within a host and scores 0.5 by construction, so this measure isolates what the window score adds beyond a host being rarely active (Section~\ref{sec:threats}). It rests on few hosts (four in the tuning scenarios) and is secondary.
\end{itemize}
Ties are broken by average rank, the expected rank under random tie-breaking; this matters for B0, under which all groups of one rule level tie. Robustness is reported with binary grades and with grades (2, 1, 1), with $\alpha = 1$, under which only the first item of each stage counts, and with each day ranked separately, since an analyst works a live queue in which an item from one day never competes with one from another.

\textbf{Accumulation (B3).} Notables are incidents and are measured as such: the time from the first attack step to the first notable containing an attack alert, the number of stages that notable spans, and on the workstation the number of notables per awake hour.

\textbf{Workstation.} The hold-out is the final week of the frozen baseline. Under stationarity 1\% of benign windows exceed the 99th baseline percentile; the observed share measures calibration drift, with a confidence interval from resampling whole days, since windows within a day are dependent. Analyst load is the number of queue items per awake hour, awake meaning outside detected downtime, whose priority exceeds a threshold fixed per system on the AIT tuning set: the priority above which half of the non-scan stages' first items fall.

\subsection{Statistical Analysis and Pre-specified Claims}
\label{sec:stats}

The test set holds about ten non-scan stage occurrences, so systems are compared \emph{paired} on the same stages: for each stage, the first-item rank under one system against the other, with the number improved and worsened and a Wilcoxon signed-rank test over the pairs. With ten pairs only a near-unanimous difference is significant, and the results are stated in those terms rather than as averaged percentages. Seeded systems are run ten times and enter the tests by their mean rank per stage. The primary comparisons are B2 against the two strongest systems without temporal context, B0$_\mathrm{host}$ and B0$_\mathrm{rank}$; all other comparisons are secondary. Holm correction is applied separately within the primary family and within the secondary family, whose comparisons are labelled exploratory; with about ten pairs the primary family then tolerates up to two losses at the smallest rank differences, where one correction over all thirteen comparisons would demand near-unanimity.

The claims are fixed in advance. Temporal context is said to improve on host or level rarity only if the primary comparison is significant after correction; otherwise the counts, the within-host AUC and the calibration split are reported and AIT-ADS is said not to separate the two. The forest is said to earn its place only if B2 beats B2$'$ after correction; otherwise the two scorers are reported as equal and the result is credited to the window representation rather than to the forest. Every paired comparison is also reported separately for stages on hosts whose calibration is dominated by their own baseline ($n_h \ge n_0$, Phase 4) and on hosts dominated by the type-level distribution: only the former test the window score as specified, the latter largely test host rarity.

\subsection{Parameter Selection}
\label{sec:params}

Table~\ref{tab:params} lists every constant of the seven phases and how its value was obtained. Values marked \emph{tuned} were selected on the tuning scenarios by the objective and veto of Section~\ref{sec:protocol}; \emph{measured} values were fixed from a measurement on data; \emph{default} values are provisional and were not varied, or were varied without changing the selected configuration.

\begin{table}[t]
\caption{Parameters of the seven phases and how each value was obtained}
\label{tab:params}
\centering\footnotesize
\setlength{\tabcolsep}{2.5pt}
\begin{tabular}{@{}c l l p{3.6cm}@{}}
\toprule
\textbf{Ph.} & \textbf{Param.} & \textbf{Value} & \textbf{Basis} \\
\midrule
3 & $N$, $\delta$ & 20, 60\,s & tuned: window sweep \\
3 & $T$ & 24\,h & tuned: 6 and 12\,h lower; 72\,h higher (0.543) but leaves one more host miscalibrated \\
3 & $G$ & 1800\,s & measured: separates 23 shutdowns \\
4 & stride & 10 & measured in simulation \\
4 & $t$, $\psi$ & 200, $\min(4096,n)$ & measured in simulation \\
4 & target / floor & 8000 / 4000 & measured in simulation (workstation) \\
4 & $n_0$ & 100 & default: 30 and 300 within one seed s.d. \\
5 & $w_r,w_c,w_s,w_n$ & 0, 0, 0.8, 0.2 & tuned: simplex, step 0.1 \\
6 & $\Delta$ & 24\,h & tuned: 15\,min to no split \\
6 & $\lambda$ (queue) & 8 & default; $\lambda = 4$ vetoed \\
7 & $t_{1/2}$ & 1\,h & tuned: 1, 4, 12\,h \\
7 & $a_r,a_c,a_s$ & 0.3, 0.2, 0.5 & tuned: ties or beats four alternatives \\
7 & $\lambda$ (accum.) & 4 & tuned: 2 and 8 give the same notables \\
7 & $\theta$ & 99.9th pct. & tuned: 99 to 99.99 give the same notables \\
7 & contrib.\ floor & 10\% & tuned: 5\% doubles benign incidents, 20\% loses notables \\
-- & $\alpha$ & 0.5 & default; $\alpha = 1$ as robustness \\
\bottomrule
\end{tabular}
\end{table}

The selection evidence for the tuned values, all on the tuning scenarios, is as follows. \emph{Hybrid weights:} $(0, 0, 0.8, 0.2)$ reaches $\alpha$-nDCG@20 0.526 against 0.501 for the sequence term alone, placing the first item of four of nine non-scan stages earlier and none later; larger novelty weights give the same ranking and the smallest is kept so that the learned term carries most of the weight; the former default $(0.3, 0.2, 0.5, 0)$ scores 0.417, and a level-percentile $S_\mathrm{rule}$ never helped. \emph{$\Delta$:} 24\,h raised $\alpha$-nDCG@20 from 0.471 (1\,h) to 0.501, with six of nine first items earlier and none later; longer gaps tie. \emph{Group score:} the maximum reaches the higher $\alpha$-nDCG@20 (0.525 against 0.501) by filling the top 20 with more stages, scans among them, but $\pi$ places eight of nine first items earlier and doubles their reciprocal rank (0.258 against 0.128), so the veto keeps $\pi$; capping it at $\lambda = 4$ (0.521) places five of nine later and is vetoed. \emph{Calibration:} shrinkage raised $\alpha$-nDCG@20 from 0.399 to 0.450 over the hard switch, and $n_0 = 30$ or 300 lies within about one seed standard deviation of $n_0 = 100$. \emph{Accumulator:} $t_{1/2} = 1$\,h raises B3 from 0.627 to 0.660 with no additional benign notable; a 5\% contribution floor gives 0.697 but doubles the benign incidents, and 20\% loses notables (0.581). \emph{Notable placement}, on four scenarios and nine stage occurrences: without promotion (a notable takes its best group's priority) B3 loses most of its gain (0.554 against 0.660), and without absorption it scores 0.641 and places five first items later and none earlier. Every choice was made before any test scenario was read.

%% ======================================================================
\section{Threats to Validity, Scope and Future Work}
\label{sec:threats}

\textbf{Threats to validity.} AIT's users and attacks are simulated, its baselines are short (no weekly cycle) and its hosts are Ubuntu servers \cite{ref22}; the workstation supplies realistic benign activity but no attacks, so no single corpus provides both. Only 6--14 non-scan attack alerts exist per scenario, so per-stage ranks are reported for every scenario rather than averaged away. The scan stages are easy for every system and are reported separately. The workstation baseline spans an examination period with reduced use, and its agent was re-enrolled once, handled by the alias table; over its five weeks the user model drifts, and periodic refitting, which a deployment would need, is not evaluated. Neither corpus has an asset inventory, so $S_\mathrm{exp}$ comes from documented host roles; it receives zero weight in the tuned priority, so no result depends on it. No attacks were run on the workstation: a positive class there would rest on one host and one attacker, the researcher, and would need its own protocol. In every AIT scenario the attack's privilege escalation lands on the intranet server, one of the least active hosts, so any signal that a host is rarely active scores well. This is a property of how AIT-ADS is built, not of attacks in general, and host rarity is therefore a comparator rather than part of the method; B0$_\mathrm{host}$, the within-host AUC and the calibration split exist to separate it from temporal context. A preliminary untuned check, not reported here, used all eight scenarios once; every later choice used the tuning scenarios only. The queue is ranked over the whole evaluation period at once, whereas an analyst works a live queue; per-day ranking is reported as a check.

\textbf{Scope and future work.} The layer consumes alerts only and reads no log content, so its reach is bounded by the rule set: four of AIT's ten attack steps never enter it. Two extensions follow. First, AIT-LDS publishes the raw logs behind every alert \cite{ref22}, so window features could include each entity's log volume and its deviation from baseline; this needs no log content and would let an unalerted step show up as unusual context around the alerts that do fire. Second, a log-anomaly detector such as AMiner could be admitted as an additional alert source whose alerts enter the same windows as Wazuh's. Neither changes the layer's role: it ranks what detectors raise and does not itself detect. A small, separately protocolled set of attacks on a monitored workstation would give the realistic corpus a positive class, and periodic refitting of the per-host models would counter the drift measured there. Robustness to quieter attacks, without reconnaissance or on busier hosts, and feedback-driven adaptation of the weights from analyst decisions \cite{ref18} are further directions.

\begin{thebibliography}{30}
\scriptsize
\setlength{\itemsep}{0pt}
\setlength{\parskip}{0pt}
\setlength{\parsep}{0pt}

\bibitem{ref1} J. Roy and E. A. Balde, ``Toward Context-Aware Alert Classification in Security Operations Centers Using LLMs,'' 2026 IEEE 5th International Conference on AI in Cybersecurity (ICAIC), Houston, TX, USA, 2026, pp. 1--10, doi: 10.1109/ICAIC67076.2026.11395735.

\bibitem{ref2} A. Hassan, A. Rauf, N. Shafqat, R. Latif, and H. Khan, ``ZenGuard a machine learning based zero trust framework for context aware threat mitigation using SIEM SOAR and UEBA,'' Scientific Reports, vol. 15, 2025, doi: 10.1038/s41598-025-20998-4.

\bibitem{ref3} A. Al Siam, M. M. Hassan, A. K. M. Masum, and T. Bhuiyan, ``Automating Malware Detection and Response via Real-Time Threat Feed Integration with Wazuh SIEM,'' 2025 IEEE 2nd International Conference on Computing, Applications and Systems (COMPAS), Kushtia, Bangladesh, 2025, pp. 1--6, doi: 10.1109/COMPAS67506.2025.11381876.

\bibitem{ref4} S. Chaurasia, A. Singh, S. K. Maurya, D. P. Bajpai, and S. Yadav, ``Graph Neural Threat Intelligence: Real-Time Attack Path Prediction Using Event Correlation from Wazuh SIEM,'' 2026 International Conference on Intelligent Processing, Hardware, Electronics, and Radio Systems (CIPHER), Jalandhar, India, 2026, pp. 1--5, doi: 10.1109/CIPHER70417.2026.11523812.

\bibitem{ref5} X. Wang, X. Yang, X. Liang, X. Zhang, W. Zhang, and X. Gong, ``Combating alert fatigue with AlertPro: Context-aware alert prioritization using reinforcement learning for multi-step attack detection,'' Computers \& Security, vol. 137, 2024, Art. no. 103583, doi: 10.1016/j.cose.2023.103583.

\bibitem{ref6} G. Suarez-Tangil, E. Palomar, A. Ribagorda, and I. Sanz, ``Providing SIEM systems with self-adaptation,'' Information Fusion, vol. 21, pp. 145--158, 2015, doi: 10.1016/j.inffus.2013.04.009.

\bibitem{ref7} M. Cinque, R. Della Corte, and A. Pecchia, ``Contextual filtering and prioritization of computer application logs for security situational awareness,'' Future Generation Computer Systems, vol. 111, pp. 668--680, 2020, doi: 10.1016/j.future.2019.09.005.

\bibitem{ref8} A. Malach, P. N. Wudali, S. Momiyama, J. Furukawa, T. Araki, Y. Elovici, and A. Shabtai, ``CyberShapley: Explanation, prioritization, and triage of cybersecurity alerts using informative graph representation,'' Computers \& Security, vol. 150, 2025, Art. no. 104270, doi: 10.1016/j.cose.2024.104270.

\bibitem{ref9} Z. Shamsi, N. V. S. Manohar, P. Likhitha, M. J. S. S. S. Jahnavi, and K. E. Rao, ``A Human Centric Risk-Based Security Alert Prioritization Using Splunk,'' 2026 IEEE Guwahati Subsection Conference (GCON), Guwahati, India, 2026, pp. 1--7, doi: 10.1109/GCON69192.2026.11649152.

\bibitem{ref10} M. Khayat, E. Barka, M. A. Serhani, F. Sallabi, K. Shuaib, and H. M. Khater, ``Advanced Techniques for Alert Management in Security Information and Event Management Systems With Ensembled Deep Learning, Hybrid Optimization, and Multi-Feature Extraction,'' IEEE Open Journal of the Communications Society, vol. 6, pp. 7349--7368, 2025, doi: 10.1109/OJCOMS.2025.3603000.

\bibitem{ref11} A. Pomsathit, ``Adaptive Detection of Advanced Persistent Threats (APT) With Graph Neural Networks and Rehearsal-Based Continual Learning on Wazuh EDR Telemetry,'' IEEE Access, vol. 13, pp. 212973--212982, 2025, doi: 10.1109/ACCESS.2025.3639270.

\bibitem{ref12} A. M. Winkler and P. Sharma, ``Proactive threat detection in enterprise systems using Wazuh: A MITRE ATT\&CK Evaluation,'' Computers \& Security, vol. 159, 2025, Art. no. 104702, doi: 10.1016/j.cose.2025.104702.

\bibitem{ref13} S. A. Chamkar, M. Zaydi, Y. Maleh, and N. Gherabi, ``Improving Threat Detection in Wazuh Using Machine Learning Techniques,'' Journal of Cybersecurity and Privacy, vol. 5, no. 2, Art. no. 34, 2025, doi: 10.3390/jcp5020034.

\bibitem{ref14} S. Ndichu, T. Ban, S. Ozawa, T. Takahashi, and D. Inoue, ``AI-Driven Security Alert Screening and Alert Fatigue Mitigation in Security Operations Centers: A Survey,'' arXiv:2605.08316 [cs.CR], May 2026, doi: 10.48550/arXiv.2605.08316.

\bibitem{ref15} T. van Ede et al., ``DEEPCASE: Semi-Supervised Contextual Analysis of Security Events,'' 2022 IEEE Symposium on Security and Privacy (SP), San Francisco, CA, USA, 2022, pp. 522--539, doi: 10.1109/SP46214.2022.9833671.

\bibitem{ref16} F. Jalalvand, M. Baruwal Chhetri, S. Nepal, and C. Paris, ``Alert Prioritisation in Security Operations Centres: A Systematic Survey on Criteria and Methods,'' ACM Computing Surveys, 2024, doi: 10.1145/3695462.

\bibitem{ref17} B. Alboushy, A. Jafar, M. Aljnidi, M. B. Disoki, and A. Shaheed, ``Context-Aware Web Attack Detection in Open-Source SIEM Systems via MITRE ATT\&CK-Enriched Behavioral Profiling,'' arXiv:2605.13337 [cs.CR], May 2026.

\bibitem{ref18} F. Jalalvand, M. Baruwal Chhetri, S. Nepal, and C. Paris, ``Adaptive alert prioritisation in security operations centres via learning to defer with human feedback,'' arXiv:2506.18462, 2025.

\bibitem{ref19} Splunk, ``How risk-based alerting works in Splunk Enterprise Security,'' Splunk Enterprise Security Documentation, accessed Sep. 2026.

\bibitem{ref20} F. T. Liu, K. M. Ting, and Z.-H. Zhou, ``Isolation Forest,'' 2008 Eighth IEEE International Conference on Data Mining, Pisa, Italy, 2008, pp. 413--422, doi: 10.1109/ICDM.2008.17.

\bibitem{ref21} M. Landauer, F. Skopik, and M. Wurzenberger, ``Introducing a New Alert Data Set for Multi-Step Attack Analysis,'' in Proc. 17th Cyber Security Experimentation and Test Workshop (CSET), 2024, doi: 10.1145/3675741.3675748. Data: AIT Alert Data Set, Zenodo, doi: 10.5281/zenodo.8263181.

\bibitem{ref22} M. Landauer, F. Skopik, M. Frank, W. Hotwagner, M. Wurzenberger, and A. Rauber, ``Maintainable Log Datasets for Evaluation of Intrusion Detection Systems,'' IEEE Transactions on Dependable and Secure Computing, vol. 20, no. 4, pp. 3466--3482, Jul.--Aug. 2023, doi: 10.1109/TDSC.2022.3201582.

\bibitem{ref23} M. Mozaffari, A. Yazdinejad, and A. Dehghantanha, ``Windows-APT 2025: A Dataset for APT-Inspired Attack Scenarios on Windows Systems,'' Data in Brief, vol. 65, Art. no. 112569, 2026, doi: 10.1016/j.dib.2026.112569.

\bibitem{ref24} C. L. A. Clarke, M. Kolla, G. V. Cormack, O. Vechtomova, A. Ashkan, S. B\"uttcher, and I. MacKinnon, ``Novelty and diversity in information retrieval evaluation,'' in Proc. 31st ACM SIGIR Conference on Research and Development in Information Retrieval, 2008, pp. 659--666, doi: 10.1145/1390334.1390446.

\bibitem{ref25} H. B\"uhlmann, ``Experience rating and credibility,'' ASTIN Bulletin, vol. 4, no. 3, pp. 199--207, 1967.

\bibitem{ref26} R. Uetz, P. B\"onninghausen, L. Hackl\"ander-Jansen, and M. Henze, ``Can Risk-Based Alerting Mitigate Cybersecurity Alert Fatigue?,'' arXiv:2609.02465 [cs.CR], Sep. 2026.

\bibitem{ref27} S. Ndichu, T. Ban, S. Ozawa, T. Takahashi, and D. Inoue, ``PACT: Reducing Alert Fatigue in Low-Prevalence SOC Streams with Triggered Active Learning,'' arXiv:2605.22324 [cs.CR], May 2026.

\bibitem{ref28} S. Freitas, J. Kalajdjieski, A. Gharib, and R. McCann, ``AI-Driven Guided Response for Security Operation Centers with Microsoft Copilot for Security,'' arXiv:2407.09017 [cs.CR], 2024.

\bibitem{ref29} S. Freitas, A. Gharib, and M. Magenheim, ``Adaptive Incident Prioritization for Security Operations at Scale,'' arXiv:2607.16963 [cs.CR], Jul. 2026.

\bibitem{ref30} N. Chalaemwongwan, ``SALAD: SOC Alert Labeled Analysis Dataset,'' King Mongkut's Institute of Technology Ladkrabang, Zenodo, Mar. 2026, doi: 10.5281/zenodo.18912306.

\bibitem{ref31} K. H. A. Hakim, ``Wazuh Alerts Dataset,'' Hugging Face Datasets, \texttt{kholil-lil/wazuh-alerts}, 2025, accessed Oct. 2026.

\end{thebibliography}
\end{document}
